% Generated draft: requires mathematical and source review.
\documentclass[11pt]{article}
\usepackage[margin=25mm]{geometry}
\usepackage{fontspec}
\setmainfont{Latin Modern Roman}
\usepackage{amsmath,amssymb,mathtools,unicode-math}
\setmathfont{Latin Modern Math}
\usepackage{enumitem,needspace,xurl,hyperref}
\setlength{\parindent}{0pt}
\setlength{\parskip}{5pt}
\setlength{\emergencystretch}{3em}
\newcommand{\N}{\mathbb{N}}
\newcommand{\Z}{\mathbb{Z}}
\newcommand{\Q}{\mathbb{Q}}
\newcommand{\R}{\mathbb{R}}
\newcommand{\C}{\mathbb{C}}
\newcommand{\F}{\mathbb{F}}
\newcommand{\A}{\mathbb{A}}
\newcommand{\PP}{\mathbb{P}}
\newcommand{\E}{\mathbb{E}}
\newcommand{\eps}{\varepsilon}
\newcommand{\dd}{\,\mathrm{d}}
\begin{document}
\section*{UnsolvedMath ID: 20003089}
\textbf{Source problem code:} catalogue entry
\subsection{Definitions and mathematical statement}
Let $S=S_{g,n}$ be an oriented finite-type punctured surface with $2g-2+n>0$, and let $\mathcal T(S)$ be its Teichm\"uller space of marked complete finite-area hyperbolic structures. For nontrivial Teichm\"uller space, the asymmetric Thurston metric is
\[
d_{\mathrm{Th}}(X,Y)=\log\sup_{\alpha}\frac{\ell_Y(\alpha)}{\ell_X(\alpha)},
\]
where $\alpha$ ranges over essential simple closed curve classes. The triply punctured sphere has a one-point Teichm\"uller space and is treated separately. The source asks whether every local isometry is induced by an extended mapping class, allowing orientation-reversing mapping classes. The original brief statement does not itself fix surface or local-isometry conventions; this is the finite-type setting of the supplied literature solution.

The aid attributes the following affirmative theorem to Pan, Theorem 1.6, and independently Huang--Ohshika--Papadopoulos, Corollary 1.17: for a connected open $U\subseteq\mathcal T(S)$, every isometric embedding $F:U\to\mathcal T(S)$ satisfying
\[
d_{\mathrm{Th}}(F(X),F(Y))=d_{\mathrm{Th}}(X,Y)
\quad(X,Y\in U)
\]
is the restriction of an element of $\operatorname{Map}^{\pm}(S)$. This preserves ordered distance, not merely a symmetrization. The theorem as reported assumes neither surjectivity onto $\mathcal T(S)$ nor differentiability and includes finite-type low-complexity cases, with the pair-of-pants case trivial. Exact admissible-surface hypotheses must be checked in the full references.

For the stronger neighborhoodwise reading of ``local'', suppose only that every $X\in U$ has an open neighborhood $V_X\subset U$ on which the displayed ordered-distance equality holds. The aid's corollary states that one effective extended mapping class still induces $F$ on all of $U$. Here the effective group is
\[
\operatorname{Map}^{\pm}(S)\big/\ker\bigl(\operatorname{Map}^{\pm}(S)\curvearrowright\mathcal T(S)\bigr).
\]
Thus $F$ automatically preserves ordered distance for every pair of points in $U$ and extends to a global mapping-class isometry, unique in this effective group. Uniqueness is not asserted for individual representatives in a possibly noneffective mapping class group. The connectedness requirement prevents independently chosen actions on separate domain components.

The extracted text continues with the historical observation attributed to Walsh that the horofunction compactification for the Thurston metric is naturally Thurston's compactification. This is retained as source context, not an extra hypothesis or the local-isometry question itself. The trailing dotted punctuation is extraction material.
\subsection{Short English statement}
The supplied literature gives an affirmative finite-type local-rigidity theorem for the ordered asymmetric Thurston metric. Its neighborhoodwise corollary yields one global mapping-class action on a connected domain, with uniqueness understood modulo the action kernel.
\subsection{Sources}
\begin{itemize}
\item Local UnsolvedMath registry, permanent ID $20003089$, source code \texttt{catalogue entry}; canonical \texttt{aim-topology-notes.json}, \texttt{notes[176]}.
\item 2014 AIM list \emph{Problems on Thurston Metric}, workshop \emph{Lipschitz metric on Teichmueller space}, Problem 1.4: \url{https://aimath.org/pastworkshops/lipschitzteichproblems.pdf}. The supplied verification reports source-PDF inspection; the registry statement status is \texttt{reconstructed\_unverified}. Original remarks and literature field are empty.
\item Pan, Theorem 1.6, and Huang--Ohshika--Papadopoulos, Corollary 1.17, as named in the aid; no titles, dates or authoritative full bibliographic entries are supplied. The extracted background separately attributes the compactification observation to Walsh.
\item Research attempt 1: \texttt{full\_solution}, status \texttt{solved}, classification \texttt{SOLVED-IN-LITERATURE}, rigor \texttt{proved}, novelty confidence low. Its literature assessment and germ-to-global corollary are machine-generated aid, not peer reviewed, and require independent verification before citation.
\end{itemize}
\end{document}
